\documentclass{article}
\usepackage{graphicx}
\usepackage{rotating}

\title{Graphics and Images Support}
\author{htex Library}
\date{2026}

\begin{document}

\maketitle

\section{Basic Image Inclusion}

You can include images using the \texttt{includegraphics} command:

\begin{figure}[h]
  \centering
  \includegraphics{example.png}
  \caption{A simple image}
  \label{fig:simple}
\end{figure}

\section{Image Sizing}

\subsection{Fixed Width}

Image with 50\% page width:

\includegraphics[width=0.5\textwidth]{image.png}

\subsection{Fixed Height}

Image with 3cm height:

\includegraphics[height=3cm]{image.png}

\subsection{Scaling}

Image scaled to 150\% of original size:

\includegraphics[scale=1.5]{image.png}

Image scaled to 0.75x:

\includegraphics[scale=0.75]{image.png}

\section{Image Rotation}

\subsection{Direct Rotation with rotatebox}

Text with rotated content:

This is \rotatebox{90}{rotated text} inline.

\subsection{Image Rotation with angle Option}

Image rotated 45 degrees:

\includegraphics[angle=45, width=5cm]{image.png}

Image rotated 90 degrees:

\includegraphics[angle=90, height=4cm]{image.png}

\section{Combined Options}

Multiple options can be combined:

\includegraphics[width=8cm, angle=270, scale=1.2]{image.png}

Width with aspect ratio preservation:

\includegraphics[width=7cm]{image.png}

\section{Figures with Captions}

\begin{figure}[h]
  \centering
  \includegraphics[width=0.8\textwidth]{photo.jpg}
  \caption{A professional photograph with caption}
  \label{fig:photo}
\end{figure}

See Figure \ref{fig:photo} on page \pageref{fig:photo}.

\begin{figure}[h]
  \centering
  \includegraphics[height=5cm, angle=180]{rotated-diagram.png}
  \caption{A diagram rotated 180 degrees}
  \label{fig:diagram}
\end{figure}

\section{Inline Images}

You can also use images inline: \includegraphics[height=0.8em]{icon.png} like this.

Images scale with text size: \large \includegraphics[height=1em]{icon.png} \normalsize

\section{Advanced Usage}

\subsection{Trimming}

Remove unwanted borders (when supported):

\includegraphics[width=6cm, trim=0 0 2cm 1cm]{image.png}

\subsection{Clipping}

Clip to specific dimensions:

\includegraphics[width=4cm, clip]{photo.jpg}

\subsection{Multiple Images in Sequence}

\includegraphics[width=0.3\textwidth]{img1.png}
\hspace{0.5cm}
\includegraphics[width=0.3\textwidth]{img2.png}
\hspace{0.5cm}
\includegraphics[width=0.3\textwidth]{img3.png}

\section{Best Practices}

\begin{enumerate}
  \item Always use meaningful file names for images
  \item Include captions in figures for clarity
  \item Use labels for cross-referencing
  \item Optimize image sizes for web presentation
  \item Ensure images are in supported formats (PNG, JPG, GIF, SVG)
\end{enumerate}

\end{document}
